% --- scholar LaTeX fragment --------------------------------------------
% Generated by: scholar sessions export -f table
% Session: t-pseudonymisation
% \input this file from the body of a document whose
% preamble loads:
%   \usepackage{longtable}
%   \usepackage{booktabs}
%   \usepackage{array}
% ----------------------------------------------------------------------

\begingroup\footnotesize
\begin{longtable}{@{}%
  >{\raggedright\arraybackslash}p{0.44\textwidth}c%
  >{\raggedright\arraybackslash}p{0.13\textwidth}%
  >{\raggedright\arraybackslash}p{0.26\textwidth}@{}}
\caption{Records bearing on the claim (3 cited, 34 supporting and 12 qualifying or contradicting it, out of 154 unique hits; the 71 adjacent and 34 false hits are not listed). The last column gives the reason for each inclusion; for machine-classified rows the model's confidence is given. Databases: OpenAlex (OA; not open access; incomplete for 3 of 3 queries), DBLP (incomplete for 3 of 3 queries), IEEE Xplore (IEEE; incomplete for 3 of 3 queries), Web of Science (WoS), Scopus, arXiv (incomplete for 3 of 3 queries).}\label{tab:t-pseudonymisation}\\
\toprule
Title & Year & Provider & Reason \\
\midrule
\endfirsthead
\multicolumn{4}{@{}l}{\emph{\tablename~\thetable{} (cont.)}}\\
\toprule
Title & Year & Provider & Reason \\
\midrule
\endhead
\midrule \multicolumn{4}{r@{}}{\emph{continued on next page}}\\
\endfoot
\bottomrule
\endlastfoot
Are `pseudonymised' data always personal data? Implications of the GDPR for administrative data research in the UK & 2018 & DBLP, IEEE, OA, arXiv & \emph{cited}: relative-approach \\
Case C-413/23 P, EDPS v. SRB: Reassessing the Role of Pseudonymization in Data Sharing? & 2026 & WoS, Scopus & \emph{cited}: relative-approach \\
The Pitfalls of Hashing for Privacy & 2018 & DBLP, IEEE, OA, arXiv & \emph{cited}: hashing-weak \\
'Potato potahto'? Disentangling de-identification, anonymisation, and pseudonymisation for health research in Africa & 2025 & WoS & supports the claim (0.80) \\
A Scalable Pseudonymization Tool for Rapid Deployment in Large Biomedical Research Networks: Development and Evaluation Study & 2024 & WoS & supports the claim (0.96) \\
A Scalable Pseudonymization Tool for Rapid Deployment in Large Biomedical Research Networks: Development and Evaluation Study & 2024 & Scopus & supports the claim (0.96) \\
Analysis of Data Anonymization Techniques & 2020 & WoS & supports the claim (0.76) \\
Anonymization and Pseudonymization & 2017 & DBLP, IEEE, OA, arXiv & supports the claim (0.88) \\
Anonymization and Pseudonymization of Personal Health Data within the Scope of Turkish Personal Data Protection Law and General Data Protection Regulation & 2024 & DBLP, IEEE, OA, arXiv & supports the claim (0.68) \\
Canadian Association of Radiologists White Paper on De-Identification of Medical Imaging: Part 1, General Principles & 2021 & WoS & supports the claim (0.70) \\
Comparative Evaluation of K-Anonymity, Differential Privacy, and Pseudonymization for Data Protection in Rare Disease Registries & 2025 & Scopus & supports the claim (0.91) \\
Comparing the benefits of pseudonymisation and anonymisation under the GDPR & 2018 & DBLP, IEEE, OA, arXiv & supports the claim (0.94) \\
Data Protection and Privacy in AI Development With a Focus on Training Data, Personal Data, Anonymisation, and Data Reuse & 2026 & Scopus & supports the claim (0.73) \\
Data Protection Using Polymorphic Pseudonymisation in a Large-Scale Parkinson's Disease Study & 2021 & DBLP, IEEE, OA, arXiv & supports the claim (0.91) \\
Data protection using strong pseudonymization & 2026 & Scopus & supports the claim (0.58) \\
Ensuring Privacy in the Labour Market: Towards Full Compliance with LGPD and GDPR & 2025 & Scopus & supports the claim (0.79) \\
Ephemeral pseudonym based de-identification system to reduce impact of inference attacks in healthcare information system & 2022 & DBLP, IEEE, OA, arXiv & supports the claim (0.92) \\
EUPID---configurable privacy-preserving record linkage in federated health data spaces & 2026 & Scopus & supports the claim (0.94) \\
GDPR-Compliant Pseudonymization Architecture for Multi-Source Healthcare Data: A Hardware-Secured Approach for Disease Registry Construction & 2026 & DBLP, IEEE, OA, arXiv & supports the claim (0.91) \\
How (not) to protect genomic data privacy in a distributed network: using trail re-identification to evaluate and design anonymity protection systems & 2004 & DBLP, IEEE, OA, arXiv & supports the claim (0.98) \\
Opening privacy sensitive microdata sets in light of GDPR & 2019 & Scopus & supports the claim (0.97) \\
Options for encoding names for data linking at the Australian Bureau of Statistics & 2018 & DBLP, IEEE, OA, arXiv & supports the claim (0.98) \\
PEEPLL: Privacy-Enhanced Event Pseudonymisation with Limited Linkability & 2019 & DBLP, IEEE, OA, arXiv & supports the claim (0.89) \\
PERSONAL DATA PROCESSING IN EDUCATIONAL INSTITUTIONS: ANONYMISATION AND PSEUDONYMISATION & 2020 & DBLP, IEEE, OA, arXiv & supports the claim (0.96) \\
Privacy Risk Evaluation of Re-Identification of Pseudonyms & 2019 & DBLP, IEEE, OA, arXiv, WoS & supports the claim (0.96) \\
Pseudonymisation and the smart city & 2018 & DBLP, IEEE, OA, arXiv & supports the claim (0.92) \\
Pseudonymisation in the context of GDPR-compliant medical research & 2023 & Scopus & supports the claim (0.93) \\
Pseudonymisation under the General Data Protection Regulation: A win-win approach? & 2017 & DBLP, IEEE, OA, arXiv & supports the claim (0.95) \\
Pseudonymization and impacts of Big (personal/anonymous) Data processing in the transition from the Directive 95/46/EC to the new EU General Data Protection Regulation & 2017 & Scopus & supports the claim (0.98) \\
Pseudonymization for research data collection: is the juice worth the squeeze? & 2019 & DBLP, IEEE, OA, arXiv & supports the claim (0.94) \\
Pseudonymization of Personal Data of Crypto Assets Users: Issues of Legal Regulation in Indonesia & 2025 & Scopus & supports the claim (0.88) \\
Pseudonymous identifiability as a societal problem & 2019 & Scopus & supports the claim (0.72) \\
Qualifying LLM tokenisation under the GDPR: pseudonymisation by default & 2026 & Scopus & supports the claim (0.95) \\
Re-pseudonymization Strategies for Smart Meter Data Are Not Robust to Deep Learning Profiling Attacks & 2024 & Scopus & supports the claim (0.99) \\
The Lawfulness of Re-identification Under Data Protection Law & 2024 & Scopus & supports the claim (0.97) \\
The Trusted Server: A secure computational environment for privacy compliant evaluations on plain personal data & 2018 & DBLP, IEEE, OA, arXiv & supports the claim (0.91) \\
What's in a name: the conflicting views of pseudonymisation under eIDAS and the General Data Protection Regulation & 2016 & DBLP, IEEE, OA, arXiv & supports the claim (0.84) \\
A quantitative approach to the GDPR's anonymisation and \textquotedbl{}appropriate technical and organisational measures\textquotedbl{} tests & 2025 & WoS & qualifies or contradicts the claim (0.96) \\
A Quantitative Approach to the GDPR's Anonymization and Pseudonymization Tests & 2025 & DBLP, IEEE, OA, arXiv & qualifies or contradicts the claim (0.88) \\
From knowing by name to targeting: the meaning of identification under the GDPR & 2022 & DBLP, IEEE, OA, arXiv & qualifies or contradicts the claim (0.94) \\
From theory to practice: Re-identification Challenge to test imaging data pseudonymization effectiveness & 2026 & Scopus & qualifies or contradicts the claim (0.94) \\
How Anonymous Is Anonymous? A Techno-Legal Exploration & 2025 & WoS & qualifies or contradicts the claim (0.95) \\
Identifiability, as a Data Risk: Is a Uniform Approach to Anonymisation About to Emerge in the EU? & 2025 & WoS & qualifies or contradicts the claim (0.98) \\
Is It Personal Data? & 2024 & DBLP, IEEE, OA, arXiv & qualifies or contradicts the claim (0.59) \\
Personal Data and Encryption in the European General Data Protection Regulation & 2016 & DBLP, IEEE, OA, arXiv & qualifies or contradicts the claim (0.95) \\
Privacy-Preserving Anonymization of System and Network Event Logs Using Salt-Based Hashing and Temporal Noise & 2025 & DBLP, IEEE, OA, arXiv & qualifies or contradicts the claim (0.88) \\
Service for the Pseudonymization of Electronic Healthcare Records Based on ISO/EN 13606 for the Secondary Use of Information & 2015 & DBLP, IEEE, OA, arXiv & qualifies or contradicts the claim (0.73) \\
The GDPR and unstructured data: is anonymization possible? & 2022 & DBLP, IEEE, OA, arXiv & qualifies or contradicts the claim (0.91) \\
The relative and contextual nature of the concept of personal data in the light of the Deloitte judgment & 2026 & Scopus & qualifies or contradicts the claim (0.97) \\
\end{longtable}
\endgroup
